\documentclass[12pt,a4paper]{article}
\usepackage{fontspec}
\XeTeXlinebreaklocale "th"
\XeTeXlinebreakskip=0pt plus 1pt
\defaultfontfeatures{Ligatures=TeX}
\setmainfont[Path={fonts/TH Sarabun PSK V-1/TH Sarabun PSK V-1/},Extension=.ttf,
  UprightFont={*},BoldFont={* Bold},ItalicFont={* Italic},BoldItalicFont={* BoldItalic}]{THSarabun}
\setmonofont{Latin Modern Mono}
\usepackage{amsmath,amssymb,graphicx,booktabs,longtable,tabularx,array}
\usepackage{geometry,fancyhdr,hyperref,xcolor,tikz,float,enumitem}
\usetikzlibrary{arrows.meta,positioning,shapes.geometric,fit}
\geometry{margin=0.82in}
\setlength{\headheight}{19pt}
\pagestyle{fancy}\fancyhf{}
\lhead{ET-FCSPS สำหรับการสลับ GFL/GFM ระดับระบบ}
\rhead{\thepage}
\hypersetup{colorlinks=true,linkcolor=blue,urlcolor=blue,citecolor=blue}
\definecolor{navy}{HTML}{17365D}
\definecolor{blue}{HTML}{2F75B5}
\definecolor{cyan}{HTML}{00A6A6}
\definecolor{green}{HTML}{2E8B57}
\definecolor{orange}{HTML}{E67E22}
\definecolor{red}{HTML}{C0392B}
\definecolor{light}{HTML}{F3F6FA}
\renewcommand{\abstractname}{บทคัดย่อ}
\renewcommand{\contentsname}{สารบัญ}
\renewcommand{\figurename}{รูปที่}
\renewcommand{\tablename}{ตารางที่}
\renewcommand{\arraystretch}{1.18}
\setlist[itemize]{leftmargin=1.45em,itemsep=.2em}
\setlist[enumerate]{leftmargin=1.65em,itemsep=.2em}
\newcommand{\currenttag}{\textcolor{green}{\textbf{VERIFIED CURRENT}}}
\newcommand{\proposedtag}{\textcolor{orange}{\textbf{\shortstack[l]{PROJECT-DERIVED\\PROPOSAL}}}}
\newcommand{\diagtag}{\textcolor{red}{\textbf{NOT YET IMPLEMENTED}}}
\newcommand{\sat}{\operatorname{sat}_{[0,1]}}
\newcommand{\R}{\mathbb{R}}

\title{\textbf{วิธีการควบคุมการสลับโหมด GFL/GFM ระดับระบบ}\\[.2em]
\large Event-Triggered Finite-Control-Set Predictive Supervisor (ET-FCSPS)\\[.3em]
\normalsize แนวทางทางเทคนิคสำหรับ IEEE 14-bus: 1 SG + 4 IBR}
\author{เอกสารข้อเสนอวิธีวิจัยและแผนพัฒนา}
\date{4 สิงหาคม 2569}

\begin{document}
\maketitle

\begin{abstract}
ระบบปัจจุบันใช้ AGSI++ เป็นดัชนีความเครียดราย IBR ร่วมกับ hysteresis และ dwell timer
ในการเปลี่ยนโหมด GFL/GFM วิธีนี้ตรวจจับความผิดปกติเฉพาะจุดได้ แต่ยังไม่ได้ตอบคำถามระดับระบบว่า
ควรเปลี่ยน IBR ตัวใด จำนวนเท่าใด และตัวใดควรเป็น angle/frequency reference leader
เพื่อให้สภาวะของโครงข่ายทั้งระบบเหมาะสมก่อนสั่งสลับจริง รายงานนี้เสนอ ET-FCSPS ซึ่งใช้ AGSI++
เป็น event indicator แล้วแจกแจง binary mode vector ของ IBR ทั้งสี่ตัว ตรวจ hard constraints
จาก full-network KCL และ reserve contract ทำนายผลของ candidate ที่ผ่านในช่วงเวลาสั้น
จากนั้นเลือกคำสั่งด้วย objective และ tie-break ที่กำหนดไว้ล่วงหน้า

เอกสารฉบับนี้อธิบายสมการ ตัวแปร ลำดับข้อมูล candidate isolation, feasibility gate,
cost function, reference transaction, fail-closed fallback, computational complexity,
pseudocode, interface และแผนทดสอบโดยละเอียด วิธีที่เสนอยังไม่ได้ implement และยังไม่มีผล
closed-loop จึงกล่าวเฉพาะประโยชน์ที่คาดหวังในรูป testable hypotheses เท่านั้น ไม่เพิ่ม AVR,
ไม่ใช้ LLM หรือ BO ในวงรอบสั่งสลับ และไม่เปลี่ยนสมการ SG/IBR ที่ผ่านการตรวจแล้ว
\end{abstract}

\tableofcontents
\newpage

\section{สถานะของเอกสารและขอบเขตหลักฐาน}

ตาราง~\ref{tab:status} แยกสิ่งที่มีหลักฐานใน repository ออกจากวิธีที่กำลังเสนอ เพื่อป้องกัน
การนำ architecture ไปกล่าวเหมือนเป็นผลทดลองแล้ว

\begin{table}[H]
\centering
\caption{สถานะขององค์ประกอบที่กล่าวถึงในรายงาน}
\label{tab:status}
\begin{tabularx}{\textwidth}{>{\bfseries}p{.27\textwidth}p{.22\textwidth}X}
\toprule
องค์ประกอบ & สถานะ & ความหมาย \\
\midrule
AGSI++ ราย IBR & \currenttag & มีสมการ bounded index, threshold, hysteresis และ dwell ใน production code \\
GFL/GFM full-state และ full KCL & \currenttag & มี route การจำลองและ transaction ที่ตรวจด้วย targeted tests \\
SG-trip override/SG handback & \currenttag & มี chronology ที่สลับ reference และ mode แบบ coordinated \\
ET-FCSPS system decision & \proposedtag & architecture, equations และ test plan ในรายงานนี้ \\
ผลว่าลด transient/จำนวน switch & \diagtag & ต้อง implement และเปรียบเทียบกับ baseline บน contract เดียวกัน \\
AVR, LLM/BO online switching & นอกขอบเขต & ไม่ถูกเพิ่มหรือใช้เป็น decision authority \\
\bottomrule
\end{tabularx}
\end{table}

คำว่า ``adaptive'' ในรายงานนี้หมายถึงการเลือก mode vector ตาม accepted system state,
topology และ reserve ณ เหตุการณ์นั้น ไม่ได้หมายถึง online parameter identification หรือการปรับ weight
โดยอัตโนมัติ ส่วนคำว่า finite-control-set ใช้หลักการประเมิน action set แบบจำกัดที่แจกแจงได้
โดยประยุกต์จากแนวคิด FCS-MPC~\cite{rodriguez2013}; ไม่ได้หมายถึงการสั่ง semiconductor switch
ระดับ converter โดยตรง

\section{ปัญหาของ AGSI++ แบบ local-only}

\subsection{สมการที่ใช้อยู่ปัจจุบัน}

สำหรับ IBR ตัวที่ $i$ production class คำนวณ sub-index ที่ PCC ดังนี้
\begin{align}
J_{V,i} &= \frac{\left||V_i|-V_{ref,i}\right|}{\Delta V_{base}}, &
J_{f,i} &= \frac{|f_i-f_0|}{\Delta f_{base}},\\
J_{R,i} &= \frac{|\dot f_i|}{\Delta R_{base}}, &
J_{P,i} &= \frac{|P_{ref,i}-P_i|}{\Delta P_{base}},\\
J_{SCR,i} &= \max\!\left(0,\frac{3}{SCR_i}-1\right), &
J_{lock,i} &= \frac{|v_{q,i}|}{\Delta v_{q,base}},\\
J_{GRA,i} &= \begin{cases}0,&GRA_i=1,\\1,&GRA_i=0.\end{cases}
\end{align}
ดัชนีรวมคือ
\begin{equation}
A_i=\sat\!\left(
w_VJ_{V,i}+w_fJ_{f,i}+w_RJ_{R,i}+w_PJ_{P,i}
+w_{SCR}J_{SCR,i}+w_{lock}J_{lock,i}+w_{GRA}J_{GRA,i}
\right),
\label{eq:agsi}
\end{equation}
โดย $\sum w=1$ และ AGSI++ preset ใช้ค่าเฉลี่ยเท่ากัน $1/7$ สำหรับเจ็ดพจน์
RoCoF ผ่าน first-order low-pass ก่อนเข้าดัชนี ส่วน raw weighted sum ถูกเก็บเพื่อวิเคราะห์แต่
decision index ถูกจำกัดให้อยู่ใน $[0,1]$ เสมอ

logic การสลับ local คือ
\begin{align}
\mathrm{GFL}\rightarrow\mathrm{GFM}:&\quad A_i\ge\Gamma_{on}
\ \text{ต่อเนื่องอย่างน้อย}\ T_{d,on},\\
\mathrm{GFM}\rightarrow\mathrm{GFL}:&\quad A_i<\Gamma_{off}
\ \text{ต่อเนื่องอย่างน้อย}\ T_{d,off},
\end{align}
โดย $\Gamma_{off}<\Gamma_{on}$ สร้าง hysteresis band ป้องกัน chatter ค่าใน report route ปัจจุบันคือ
$\Gamma_{on}=0.65$, $\Gamma_{off}=0.35$, $T_{d,on}=0.10$~s และ $T_{d,off}=1.00$~s
ค่าดังกล่าวจะถูก freeze เป็น baseline; รายงานนี้ไม่อนุญาตให้เปลี่ยนเพื่อทำให้ผลของวิธีใหม่ดีขึ้น

\subsection{AGSI++ มองเห็นอะไรและยังไม่เห็นอะไร}

AGSI++ ไม่ใช่ข้อมูลน้อย: มันมี voltage, frequency, RoCoF, active-power tracking, SCR,
PLL alignment และ reference availability แต่ทุกพจน์ผูกกับ IBR ตัวเดียวหรือ PCC ของตัวนั้น
จึงยังไม่มีการประเมินพร้อมกันว่า
\begin{itemize}
  \item voltage violation เกิดที่บัสอื่นหรือมีขอบเขตกว้างเพียงใด;
  \item topology ปัจจุบันทำให้ IBR ตัวใดอยู่ใกล้พื้นที่อ่อนในความหมายทางไฟฟ้า;
  \item active/reactive reserve ของชุด IBR ที่จะเปลี่ยนโหมดเพียงพอหรือไม่;
  \item การสลับหนึ่งตัวหรือหลายตัวจะเพิ่ม current-limit exposure ที่จุดอื่นหรือไม่;
  \item energized island ทุกเกาะมี reference owner และ owner ที่เลือกมี headroom หรือไม่;
  \item candidate หนึ่งจะให้ transient หลัง commit ดีกว่า candidate อื่นหรือไม่
\end{itemize}
ดังนั้นข้อจำกัดไม่ใช่ ``AGSI ไม่มีแรงดันหรือกำลัง'' แต่คือมันยังไม่ได้ใช้ full-network consequence
เป็น decision authority ก่อน commit

\section{แนวคิดหลักของ ET-FCSPS}

ET-FCSPS วางชั้นตัดสินใจระดับระบบเหนือ local state machine โดยไม่ยกเลิก AGSI++
โครงสร้างเป็นไปตามรูป~\ref{fig:architecture}

\begin{figure}[H]
\centering
\begin{tikzpicture}[>=Latex,node distance=8mm,
box/.style={rounded corners=3pt,draw=blue,thick,fill=light,minimum width=10.4cm,
minimum height=.85cm,align=center},
side/.style={rounded corners=3pt,draw=orange,thick,fill=orange!8,minimum width=3.1cm,
minimum height=.8cm,align=center}]
\node[box] (meas) {accepted production state: $x_k,y_k$, topology, modes, timers, limits, reserve metadata};
\node[box,below=of meas] (local) {AGSI++ ราย IBR + scheduled/topology/reference events $\Rightarrow$ event request};
\node[box,below=of local] (sup) {ET-FCSPS: enumerate $\rightarrow$ prune $\rightarrow$ KCL screen $\rightarrow$ predict $\rightarrow$ rank};
\node[box,below=of sup] (txn) {hybrid transaction: state mapping + right-limit KCL + dwell/guard + atomic commit};
\node[side,right=10mm of sup] (fail) {ไม่มี feasible candidate\\fail-closed fallback};
\draw[->,thick] (meas)--(local);
\draw[->,thick] (local)--(sup);
\draw[->,thick] (sup)--(txn);
\draw[->,thick,red] (sup)--(fail);
\end{tikzpicture}
\caption{สถาปัตยกรรมที่เสนอ: AGSI++ เป็น local evidence ส่วน ET-FCSPS เป็นผู้อนุมัติคำสั่งระดับระบบ}
\label{fig:architecture}
\end{figure}

หลักสำคัญมีสี่ข้อ
\begin{enumerate}
  \item \textbf{Event-triggered}: ไม่รัน candidate search ทุก integration step แต่รันเมื่อมี local request,
  topology/reference event หรือ post-event re-evaluation ตาม contract
  \item \textbf{Finite action set}: แจกแจง mode vector ที่เป็นไปได้ทั้งหมด จึง audit ได้ว่าได้ตรวจตัวเลือกใด
  \item \textbf{Constraint first}: candidate ที่ผิด KCL, limit, reserve หรือ reference gate ถูกตัดก่อนดู cost
  \item \textbf{Predict before commit}: candidate trial ทำงานบนสำเนา state; state จริงเปลี่ยนเฉพาะ atomic commit
\end{enumerate}

\section{นิยามข้อมูลเข้าและ system snapshot}

เมื่อ controller ถูกเรียกที่ accepted sample $k$ ให้สร้าง snapshot
\begin{equation}
\mathcal Z_k=\{t_k,x_k,y_k,\mathcal T_k,\mathbf m_k,\boldsymbol\tau_k,
\mathcal L_k,\mathcal R_k,\mathcal E_k\},
\end{equation}
โดยตัวแปรมีความหมายตามตาราง~\ref{tab:snapshot}

\begin{longtable}{>{\bfseries}p{.17\textwidth}p{.76\textwidth}}
\caption{ข้อมูลใน accepted system snapshot}\label{tab:snapshot}\\
\toprule
สัญลักษณ์ & ความหมาย \\
\midrule\endfirsthead
\toprule สัญลักษณ์ & ความหมาย \\
\midrule\endhead
$t_k$ & เวลาของ accepted sample; ห้ามใช้ rejected parent trial \\
$x_k$ & differential states ของ SG และ IBR ตาม production state order \\
$y_k$ & algebraic network voltages/currents ของ full-KCL solve \\
$\mathcal T_k$ & topology, energized islands, line/breaker status และ event side \\
$\mathbf m_k$ & committed mode vector ของ IBR ณ ปัจจุบัน \\
$\boldsymbol\tau_k$ & up/down dwell timers, lockout และ last-switch time ราย IBR \\
$\mathcal L_k$ & voltage/current/operating limits ที่มี provenance \\
$\mathcal R_k$ & active/reactive reserve และ current headroom; ถ้าไม่ทราบต้องเป็น unknown ไม่ใช่ $\infty$ \\
$\mathcal E_k$ & event identity, AGSI parts, GRA และ reference owner ปัจจุบัน \\
\bottomrule
\end{longtable}

global metrics ที่ controller ใช้ต้องคำนวณจาก base และ sign convention เดียวกับ KCL เช่น
\begin{align}
V_{min,k}&=\min_b |V_{b,k}|,& V_{max,k}&=\max_b |V_{b,k}|,\\
\Delta P_{avail,k}&=P_{SG}^{avail}+\sum_iP_i^{avail}-P_{load}-P_{loss},\\
\Delta Q_{avail,k}&=Q_{SG}^{avail}+\sum_iQ_i^{avail}-Q_{load}-Q_{loss}.
\end{align}
ต้องไม่ใช้ $\sum P_{electrical}-P_{load}-P_{loss}$ จาก algebraic solution เป็น reserve metric เพราะ KCL
ทำให้ผลรวมนั้นใกล้ศูนย์โดยนิยาม สิ่งที่ต้องใช้คือ available/commandable source power และ headroom
ซึ่งต้องมี source/case contract หาก energy-source model ยังไม่ระบุ candidate ที่ต้องพึ่ง reserve ดังกล่าว
ต้องถูกจัดเป็น ``ไม่ทราบ'' และไม่อาจผ่าน production gate ด้วยการสมมติค่าเอง

\section{decision variables และ candidate universe}

สำหรับ IBR สี่ตัว กำหนด
\begin{equation}
\mathbf m=[m_1,m_2,m_3,m_4]\in\{0,1\}^4,
\qquad m_i=0:\mathrm{GFL},\quad m_i=1:\mathrm{GFM}.
\end{equation}
จึงมี mode vector $2^4=16$ รูปแบบ ถ้า SG online และเป็น reference owner คงที่ candidate universe
มีไม่เกิน 16 รูปแบบ ถ้า SG offline และยอมให้ GFM ทุกตัวในชุดเป็น owner candidate ได้ จำนวนคู่
$(\mathbf m,r)$ สูงสุดคือ
\begin{equation}
\sum_{\emptyset\ne S\subseteq\{1,2,3,4\}}|S|
=4\,2^{4-1}=32.
\end{equation}
นี่เป็นขนาดเล็กพอให้แจกแจงครบโดยไม่ต้องใช้ black-box optimiser อย่างไรก็ตาม candidate ที่มี
mode เดิม, ติด dwell/lockout, ไม่มี owner หรือไม่มี state mapping ที่ valid จะถูก prune ก่อน prediction

reference variable นิยามเป็น
\begin{equation}
r\in\{\mathrm{SG},1,2,3,4\}.
\end{equation}
ถ้า $r=i$ ต้องมี $m_i=1$ และ energized island นั้นต้องมี owner เพียงหนึ่งตัว คำว่า owner หมายถึง
ผู้ให้ angle/frequency reference ใน formulation ของ island ไม่ใช่การเปลี่ยน IBR ให้เป็น algebraic PF slack;
ทุกอุปกรณ์ยังอยู่ใน full-network KCL

\section{event trigger และ candidate generation}

supervisor ถูกเรียกเมื่ออย่างน้อยหนึ่งเงื่อนไขเป็นจริง
\begin{enumerate}
  \item local state machine รายงาน pending GFL$\rightarrow$GFM หรือ GFM$\rightarrow$GFL request;
  \item SG/breaker/reference availability เปลี่ยน;
  \item topology หรือ load event ถูก apply/restore;
  \item fault apply/clear;
  \item synchronism guard ครบ dwell และต้องประเมิน handback;
  \item post-commit verification interval ถึงเวลาตาม contract
\end{enumerate}
ไม่มีการสร้าง threshold ใหม่เพื่อเรียก controller หากยังไม่มี source/approved derivation; ในเฟสแรกให้ reuse
event identities และ local requests ที่มีอยู่

candidate generation ทำตามลำดับ
\begin{enumerate}
  \item ระบุ energized islands จาก $\mathcal T_k$
  \item แจกแจง $\mathbf m$ และ owner $r$
  \item ตัด candidate ที่ขัดกับ device capability, dwell, lockout และ reference rule
  \item map state บนสำเนา $x_k$ ตาม GFL/GFM transfer contract
  \item สร้าง candidate fingerprint จาก snapshot identity, mode vector, owner และ option contract
\end{enumerate}

\section{hard feasibility gates}

candidate $c=(\mathbf m,r)$ ต้องผ่านทุก gate ก่อนคำนวณ soft cost
\begin{align}
g_{finite}(c)&:\quad x^{(c)},y^{(c)}\ \text{finite},\\
g_{KCL}(c)&:\quad \|F(x^{(c)},y^{(c)},c)\|_\infty\le\varepsilon_{KCL},\\
g_V(c)&:\quad V_{min}^{lim}\le |V_b^{(c)}|\le V_{max}^{lim}\quad\forall b,\\
g_I(c)&:\quad |I_i^{(c)}|\le I_{i,max}\quad\forall i,\\
g_{ref}(c)&:\quad \text{ทุก energized island มี authenticated owner หนึ่งตัว},\\
g_{reserve}(c)&:\quad \Delta P_{avail}^{(c)}\ge0,\quad\Delta Q_{avail}^{(c)}\ge0,\\
g_{hyb}(c)&:\quad \text{state mapping, dwell, lockout และ transaction identity valid}.
\end{align}
ค่าขอบเขตและ tolerance ต้องมาจาก production/case contract และถูก freeze ก่อนเห็นผล comparison
ห้ามเลือก candidate ที่ผิด hard constraint เพียงเพราะ cost ต่ำกว่า

\subsection{algebraic screening}

ชั้นแรกใช้ project-owned PF/full-KCL solver ประเมิน right-limit ของ candidate โดยไม่ advance เวลา
ผลลัพธ์ต้องเก็บ converged flag, residual, voltage range, current margin, reserve status และ rejection reason
candidate ที่ Newton ไม่ converge ต้องถูก reject ไม่ใช่แทนด้วยค่าจาก candidate ก่อนหน้า

\subsection{สิ่งที่ยังไม่ใช่ feasibility proof}

PF convergence อย่างเดียวไม่พิสูจน์ transient stability และ local eigenvalue อย่างเดียวไม่พิสูจน์ผลของ
multi-IBR island ดังนั้น candidate ที่ผ่าน algebraic gate ยังต้องผ่าน short-horizon prediction และ runtime
domain gates การที่ objective ต่ำก็ไม่ใช่ข้อยกเว้นให้ข้ามขั้นเหล่านี้

\section{short-horizon prediction}

สำหรับ candidate ที่ผ่าน algebraic screen ให้สร้าง isolated trial context
\begin{equation}
(\hat x_0^{(c)},\hat y_0^{(c)})=\mathcal M_c(x_k,y_k),
\end{equation}
โดย $\mathcal M_c$ คือ audited mode-state mapping แล้วจำลอง
\begin{align}
\dot{\hat x}^{(c)} &= f(\hat x^{(c)},\hat y^{(c)},c,u_k),\\
0 &= F(\hat x^{(c)},\hat y^{(c)},c,\mathcal T_k),
\qquad t\in[t_k,t_k+T_p].
\end{align}
$T_p$ เป็น prediction horizon ที่ต้องประกาศก่อน experiment ในเฟสแรกให้ทดสอบหลายค่าที่ predeclared
เพื่อรายงาน sensitivity ไม่เลือกค่าหลังเห็นว่า candidate ใดชนะ

กฎ isolation คือ
\begin{itemize}
  \item trial ทุกตัวเริ่มจาก accepted snapshot เดียวกัน;
  \item ห้ามแก้ production device object, timer, RoCoF memory หรือ event log ระหว่าง trial;
  \item trial ที่ออกจาก validity domain, residual gate หรือ limit ต้องหยุดและบันทึก fail reason;
  \item หลังประเมินแล้วทิ้ง trial state ทั้งหมด; commit ได้เพียง winner หลังตรวจ fingerprint ซ้ำ
\end{itemize}

\section{objective function และ normalization}

ให้ $\mathcal C_k^{feas}$ เป็นชุด candidate ที่ผ่าน hard gates เลือก
\begin{equation}
c_k^*=\arg\min_{c\in\mathcal C_k^{feas}}J(c;\mathcal Z_k),
\end{equation}
โดยตัวอย่าง objective ที่เสนอคือ
\begin{align}
J(c)=&\;w_VJ_V(c)+w_fJ_f(c)+w_RJ_{RoCoF}(c)+w_PJ_{P,res}(c)
+w_QJ_{Q,res}(c)\notag\\
&+w_IJ_I(c)+w_HJ_{hand}(c)+w_NJ_{GFM}(c)+w_SJ_{switch}(c).
\label{eq:cost}
\end{align}
พจน์ทั้งหมดต้อง dimensionless ก่อนถ่วงน้ำหนัก ตัวอย่างนิยามคือ
\begin{align}
J_V(c)&=\max_{t,b}\left[
\frac{(|V_b|-V_{max}^{lim})_+ +(V_{min}^{lim}-|V_b|)_+}{\Delta V_{norm}}
\right],\\
J_f(c)&=\max_t\frac{|f_{ref}(t)-f_0|}{\Delta f_{norm}},\\
J_{RoCoF}(c)&=\max_t\frac{|\dot f_{ref}(t)|}{\Delta R_{norm}},\\
J_I(c)&=\max_{t,i}\frac{|I_i(t)|}{I_{i,max}},\\
J_{GFM}(c)&=\frac{1}{4}\sum_i m_i,\\
J_{switch}(c)&=\frac{1}{4}\sum_i|m_i-m_{i,k}|,\\
J_{hand}(c)&=\max\left(
\frac{|\Delta P|}{\Delta P_{norm}},
\frac{|\Delta Q|}{\Delta Q_{norm}},
\frac{|\Delta I|}{\Delta I_{norm}},
\frac{|\Delta\theta|}{\Delta\theta_{norm}}
\right).
\end{align}
เมื่อ $(a)_+=\max(a,0)$ ส่วน $J_{P,res}$ และ $J_{Q,res}$ วัดระยะใกล้ขอบ reserve ไม่ใช่
electrical KCL mismatch

\subsection{กติกาการกำหนด weight}

weight ทุกตัวต้องไม่ติดลบและ $\sum w=1$ การกำหนดต้องจัดประเภท SOURCE\_DEFINED,
CASE\_DEFINED หรือ PROJECT\_DERIVED และ freeze ก่อนผล comparison หากยังไม่มีเหตุผลรองรับ
ให้รายงาน sensitivity grid ที่กำหนดล่วงหน้าแทนการหา weight ที่ทำให้ผลสวยที่สุด
BO อาจใช้เป็น offline comparison ได้ในอนาคต แต่ไม่ใช่ส่วนจำเป็นของ ET-FCSPS และ output ของ BO
ห้าม feed production ก่อนผ่าน provenance และ validation contract

\section{deterministic ranking และ tie-break}

เพื่อให้ทำซ้ำได้ candidate ordering ใช้ลำดับต่อไปนี้
\begin{enumerate}
  \item ผ่าน hard gates ทุกข้อก่อน candidate ที่ไม่ผ่าน;
  \item $J(c)$ ต่ำกว่า;
  \item จำนวน switching actions ต่ำกว่า;
  \item จำนวน GFM ต่ำกว่า;
  \item reserve margin ต่ำสุดของ candidate สูงกว่า;
  \item accepted residual ต่ำกว่า;
  \item lexicographic mode vector และ owner ID เป็น tie-break สุดท้าย
\end{enumerate}
ต้องกำหนด floating-point comparison tolerance ล่วงหน้าและเก็บ ordering key ของทุก candidate
ห้ามพึ่ง iteration order ของ cell/struct หรือ random seed

\section{reference leader และ transaction semantics}

\subsection{SG trip หรือ reference loss}

เมื่อ SG เปิดและ $GRA=0$ controller ต้องเลือก $(\mathbf m,r)$ ที่มีอย่างน้อยหนึ่ง GFM
จากนั้น commit เป็น transaction เดียว
\begin{enumerate}
  \item freeze accepted left-limit state;
  \item authenticate event และ candidate fingerprint;
  \item map ทุก branch ตาม mode vector;
  \item กำหนด owner $r$ และ reference frame;
  \item solve right-limit KCL;
  \item ถ้าผ่านทุก gate จึงเปลี่ยน committed state/mode/timer พร้อมกัน;
  \item ถ้าไม่ผ่าน rollback ทั้ง transaction
\end{enumerate}

\subsection{SG reclose และ handback}

การคืน reference ต้องตรวจ synchronism guard
\begin{equation}
|\Delta V|\le\Delta V_{max},\qquad
|\Delta f|\le\Delta f_{max},\qquad
|\Delta\theta|\le\Delta\theta_{max}
\end{equation}
ครบ dwell ก่อน close จากนั้น controller ประเมิน candidate ที่ SG เป็น owner ซึ่งอาจเป็น all-GFL
หรือคงบาง IBR เป็น GFM ชั่วคราว หาก policy อนุญาต การ handback ต้องใช้ state mapping และ right-limit
KCL เดียวกับ forward switch เพื่อลด jump และป้องกัน partial commit

\section{fail-closed fallback}

ถ้า $\mathcal C_k^{feas}=\emptyset$ ห้ามเลือก candidate ที่ผิด constraint แผน fallback เฟสแรกคือ
\begin{enumerate}
  \item ถ้า current committed configuration ยังผ่าน hold gate ให้คง mode เดิมและบันทึก no-decision;
  \item เมื่อเกิด authenticated SG/reference loss ให้ลอง current verified atomic all-GFM emergency route
  เฉพาะเมื่อ route นั้นผ่าน right-limit gates;
  \item ถ้า emergency route ไม่ผ่าน ให้ simulation/controller return explicit infeasible status;
  \item ห้ามเติม trajectory, freeze graph หรือแทน state ด้วย candidate อื่น;
  \item protection action ของระบบจริงอยู่นอก model และต้องออกแบบแยกต่างหาก
\end{enumerate}

\section{pseudocode ของ controller}

\begin{figure}[H]
\small
\begin{minipage}{.94\textwidth}
\ttfamily
01  function decision = et\_fcs\_supervisor(accepted\_snapshot, event)\par
02  assert(snapshot is accepted, finite, fingerprinted)\par
03  if not local\_request and not authenticated\_event: return HOLD\par
04  C = enumerate\_mode\_owner\_pairs(snapshot)\par
05  C = prune\_by\_capability\_dwell\_lockout\_reference(C)\par
06  feasible = []\par
07  for each c in C in canonical order\par
08      trial = deep\_copy(snapshot)\par
09      trial = map\_states\_for\_candidate(trial,c)\par
10      a = solve\_right\_limit\_full\_KCL(trial,c)\par
11      if not algebraic\_gates\_pass(a): log reject; continue\par
12      p = simulate\_short\_horizon\_isolated(a,c,Tp)\par
13      if not dynamic\_gates\_pass(p): log reject; continue\par
14      feasible.append(c, metrics(p), cost(p), ordering\_key)\par
15  end\par
16  if feasible is empty: return declared\_fail\_closed\_fallback()\par
17  winner = canonical\_argmin(feasible)\par
18  revalidate(snapshot\_fingerprint, event, winner)\par
19  return atomic\_commit\_request(winner, evidence\_table)\par
20  end\par
\end{minipage}
\caption{pseudocode ที่เสนอ; ชื่อฟังก์ชันเป็น interface plan ยังไม่ใช่ไฟล์ production}
\end{figure}

\section{computational complexity และการลด runtime}

ถ้าประเมินครบ $N_c$ candidates และแต่ละ prediction ใช้ $N_p=T_p/\Delta t$ steps ต้นทุนโดยคร่าวคือ
\begin{equation}
\mathcal O\!\left(N_c\,[C_{KCL}+N_pC_{step}]\right),
\end{equation}
โดย $N_c\le16$ เมื่อ SG เป็น owner และ $N_c\le32$ เมื่อเลือก IBR owner ด้วย วิธีลดเวลาโดยไม่ลด
numerical integrity คือ
\begin{enumerate}
  \item logical pruning ก่อนเรียก solver;
  \item algebraic screen ก่อน dynamic prediction;
  \item หยุด candidate ทันทีเมื่อผิด hard gate;
  \item ทำนายเฉพาะ event-trigger ไม่ใช่ทุก time step;
  \item parallel candidate evaluation ได้เมื่อ trial object แยกกันจริงและผล sorting deterministic;
  \item cache ได้เฉพาะ immutable equilibrium ที่ fingerprint ตรงทุก input
\end{enumerate}
production chronology ปัจจุบันใช้เวลาระดับหลายสิบนาที จึงยังห้ามอ้าง real-time capability
ก่อนมี profiling ของ controller deadline, worst-case candidate count และ solver iteration distribution

\section{ผลลัพธ์ที่คาดหวังและวิธีพิสูจน์}

ผลที่คาดหวังในตาราง~\ref{tab:expected} เป็น hypotheses ไม่ใช่ผลทดลอง

\begin{longtable}{>{\bfseries}p{.23\textwidth}p{.37\textwidth}p{.32\textwidth}}
\caption{ประโยชน์ที่คาดหวัง กลไก และ metric}\label{tab:expected}\\
\toprule
ผลที่คาดหวัง & กลไก & metric \\
\midrule\endfirsthead
\toprule ผลที่คาดหวัง & กลไก & metric \\
\midrule\endhead
ลดการสลับที่ไม่จำเป็น & objective คำนึงถึงจำนวน GFM และจำนวนครั้งที่เปลี่ยน mode & switch count, chatter, GFM-on time \\
ลด voltage/frequency transient & ทำนายก่อน commit และตัด candidate ที่ excursion สูง & peak $|\Delta f|$, RoCoF, voltage violation, settling time \\
ลด current/limiter stress & current/reserve เป็น hard gate และ cost margin & $\max |I|/I_{max}$, limiter dwell, reserve margin \\
ทำ handback ให้นุ่มขึ้น & เปรียบเทียบ jump ของ candidate ก่อน atomic commit & $\Delta P,\Delta Q,\Delta I,\Delta\theta$ ที่ event \\
รักษา reference continuity & บังคับทุก energized island มี owner & no-reference duration ต้องเป็นศูนย์ \\
อธิบายย้อนหลังได้ & เก็บ candidate evidence และ deterministic ordering & decision table reproduce ได้จาก snapshot \\
\bottomrule
\end{longtable}

การพิสูจน์ต้องใช้ baseline AGSI-only และ ET-FCSPS บน case, event, parameter, $\Delta t$,
solver options และ gates เดียวกัน รายงานทั้ง improvement และ regression ห้ามเลือกเฉพาะเหตุการณ์ที่ method ใหม่ชนะ

\section{การเปรียบเทียบกับ Bayesian optimisation}

BO สร้าง probabilistic surrogate ของ expensive black-box objective แล้วใช้ acquisition function เลือกจุดทดลอง
ถัดไป~\cite{frazier2018,snoek2012} จึงเหมาะกับ offline parameter/hyperparameter search ที่ evaluation แพง
ขณะที่ constrained MPC ใช้ current state เป็น initial condition แก้ finite-horizon problem และใช้ action แรก
โดยมีจุดเด่นด้าน state/control constraints~\cite{mayne2000}

\begin{table}[H]
\centering
\caption{บทบาทของ BO และ ET-FCSPS ในโจทย์นี้}
\begin{tabularx}{\textwidth}{>{\bfseries}p{.22\textwidth}XX}
\toprule
ประเด็น & BO & ET-FCSPS \\
\midrule
โจทย์หลัก & ค้น parameter ของ expensive black-box objective & เลือก action จาก finite mode/owner set ณ accepted state \\
ข้อมูล & evaluations เพื่อสร้าง surrogate/uncertainty & project plant model + full-network snapshot \\
constraint & ต้องออกแบบ constrained BO เพิ่ม & reject hard-constraint violation ก่อน ranking \\
traceability & ต้องอธิบาย surrogate, kernel และ acquisition & แสดง gate/metric/cost ของทุก candidate \\
ต้นทุนหลัก & จำนวน expensive samples & จำนวน candidate คูณ prediction horizon \\
บทบาทที่เหมาะ & offline design/tuning study & event-time supervisory action selection \\
\bottomrule
\end{tabularx}
\end{table}

จึงควรอธิบายว่า ET-FCSPS ``เหมาะกับ decision set เล็กและต้อง audit ได้ในงานนี้'' ไม่ใช่ ``ดีกว่า BO
ทุกกรณี'' หากงานรุ่นก่อนใช้ BO เพื่อปรับ weight offline ทั้งสองวิธีกำลังทำคนละหน้าที่ และ comparison ที่เป็นธรรม
ต้องใช้ objective/data budget/event set เดียวกัน

\section{แผน interface สำหรับ implementation}

ชื่อไฟล์ในตาราง~\ref{tab:interfaces} เป็นข้อเสนอ ยังไม่มี production implementation

\begin{table}[H]
\centering
\caption{interface ที่เสนอและหน้าที่}
\label{tab:interfaces}
\begin{tabularx}{\textwidth}{>{\raggedright\arraybackslash}p{.42\textwidth}X}
\toprule
ชื่อที่เสนอ & หน้าที่ \\
\midrule
\texttt{+stability/et\_fcs\_snapshot.m} & สร้าง immutable accepted snapshot และ fingerprint \\
\texttt{+stability/et\_fcs\_enumerate.m} & แจกแจง mode/owner pairs และ canonical IDs \\
\texttt{+stability/et\_fcs\_screen.m} & state mapping, right-limit KCL และ hard gates \\
\texttt{+stability/et\_fcs\_predict.m} & isolated short-horizon trial \\
\texttt{+stability/et\_fcs\_metrics.m} & global voltage/frequency/current/reserve metrics \\
\texttt{+stability/et\_fcs\_rank.m} & normalization, cost และ deterministic tie-break \\
\texttt{+stability/et\_fcs\_supervisor.m} & orchestration และ fail-closed decision \\
\texttt{tests/test\_et\_fcs\_*.m} & unit/oracle/integration/falsification tests \\
\bottomrule
\end{tabularx}
\end{table}

output schema ขั้นต่ำต้องมี event ID, snapshot fingerprint, candidate ID, mode vector, owner,
prune/reject reason, right-limit residual, predicted metrics, normalized costs, ordering key,
winner identity, runtime และ commit/rollback result

\section{แผนพัฒนาเป็นเฟส}

\begin{enumerate}
  \item \textbf{Freeze contract}: state order, base/sign, limits, reserve semantics, event identity และ baseline outputs
  \item \textbf{Global metrics}: implement voltage/topology/reference/current metrics พร้อม independent identities
  \item \textbf{Enumeration oracle}: พิสูจน์ 16 mode vectors และ 32 mode-owner pairs ไม่ซ้ำและ ordering deterministic
  \item \textbf{Algebraic screening}: clone/mapping/right-limit KCL โดยไม่มี side effect
  \item \textbf{Prediction and cost}: short-horizon isolation, metric normalization และ predeclared weight set
  \item \textbf{Transaction integration}: atomic commit, fingerprint revalidation, rollback และ fallback
  \item \textbf{Falsification}: no-event, load, fault, line trip, SG trip/reclose, dt/horizon/runtime sensitivity
  \item \textbf{Baseline comparison}: AGSI-only vs ET-FCSPS บน identical contracts
\end{enumerate}

\section{verification matrix}

\begin{longtable}{>{\bfseries}p{.20\textwidth}p{.39\textwidth}p{.33\textwidth}}
\caption{เกณฑ์ตรวจที่ต้องประกาศก่อนผล}\label{tab:verification}\\
\toprule
Gate & สิ่งที่ตรวจ & หลักฐาน \\
\midrule\endfirsthead
\toprule Gate & สิ่งที่ตรวจ & หลักฐาน \\
\midrule\endhead
Enumeration & mode/owner universe ครบ ไม่ซ้ำ & independent combinatorial oracle \\
Snapshot isolation & trial ไม่เปลี่ยน production object/state/timer & before/after fingerprint bit identity \\
Metric identity & $P,Q,loss,current,reserve$ ใช้ base/sign ถูกต้อง & direct KCL/power identities \\
Feasibility & candidate ผิด constraint ถูก reject & injected counterexamples ทุก gate \\
Ranking & winner เป็น canonical minimum & independent sorted table oracle \\
Hybrid transaction & mapping/right-limit/rollback atomic & no partial commit และ finite residual \\
Reference security & ทุก energized island มี owner & zero no-reference accepted duration \\
Behaviour & no-event/load/fault/line/SG events & ไม่มี false switch/chatter และ log มีเหตุผล \\
Numerical & $\Delta t,T_p$, finite-difference sensitivity & รายงานความต่างโดยไม่ tune tolerance \\
Runtime & candidate count, solver iterations, deadline & percentile/worst-case พร้อม timeout fallback \\
Comparison & baseline และ proposal ใช้ contract เดียวกัน & paired scenario table ไม่คัดเฉพาะผลดี \\
\bottomrule
\end{longtable}

full repository regression ไม่จำเป็นทุก edit ระหว่างพัฒนา แต่ targeted gates ต้องครอบคลุม producer,
consumer และ failure path ทุกเฟส หลัง integration ที่เปลี่ยน shared runtime ให้ใช้ risk policy ของ repository
กำหนดว่าต้องขยาย regression หรือไม่

\section{ตัวอย่างลำดับการตัดสินใจ}

\subsection{SG trip}

ที่ event SG trip, $GRA$ เปลี่ยนเป็นศูนย์และ controller ถูก trigger สมมติ logical pruning เหลือ candidate
ที่มี GFM อย่างน้อยหนึ่งตัวและ owner อยู่ใน GFM set จากนั้น KCL screen อาจตัดชุดที่ voltage/current gate
ไม่ผ่าน prediction เปรียบเทียบชุดที่เหลือ หาก $(1000,r=1)$ ผ่านและ cost ต่ำที่สุด controller จึง commit
IBR1 เป็น GFM owner ส่วน IBR2--4 คง GFL ไม่ใช่สั่ง all-GFM อัตโนมัติ แต่ถ้าชุดหนึ่งตัวไม่ผ่านและ
$(1111,r=1)$ เป็นชุดเดียวที่ feasible การเลือก all-GFM ก็มีเหตุผลจาก evidence table ไม่ใช่เพราะ AGSI ชนพร้อมกัน

\subsection{fault เฉพาะจุด}

AGSI ของ IBR ใกล้ fault อาจ trigger ก่อนตัวอื่น แต่ supervisor ยังประเมิน bus-voltage violation,
electrical position, current headroom และ predicted transient ของทุก admissible set ผลอาจเลือก IBR ใกล้ fault,
เลือก IBR ที่มี reserve มากกว่า หรือไม่สลับเลยถ้า ride-through configuration ผ่านทุก gate ทั้งหมดเป็น hypothesis
จนกว่าจะรัน predictor จริง

\subsection{SG reclose}

เมื่อ synchronism guard ผ่าน supervisor สร้าง candidate ที่ SG กลับเป็น owner เปรียบเทียบ all-GFL กับชุดที่
คง GFM ชั่วคราวโดยพิจารณา handback jump และ transient ถ้า all-GFL ผ่านและมี cost ต่ำจึง commit คืนโหมด
หากไม่มี candidate ผ่านต้อง hold/fail-closed ไม่บังคับกลับเพียงเพราะ $A_i<\Gamma_{off}$

\section{แนวตอบคำถามทางเทคนิค}

\subsection*{ถาม: ยังใช้ AGSI++ อยู่หรือไม่?}
ใช้ แต่เปลี่ยนบทบาทจากผู้สั่งโดยตรงเป็น local event indicator และ evidence ราย IBR การอนุมัติคำสั่งอยู่ที่
system-level supervisor

\subsection*{ถาม: adaptive ตรงไหนถ้า weight ไม่เปลี่ยน?}
คำสั่ง mode/owner ปรับตาม accepted network state, topology, limits และ predicted consequences ทุก event
จึงเป็น state-adaptive action selection ไม่ใช่ adaptive parameter identification

\subsection*{ถาม: ทำไมไม่ใช้ LLM?}
คำสั่งต้อง deterministic, bounded, replayable และ fail-closed LLM ไม่มี numerical/constraint guarantee
จึงเหมาะกับการอธิบาย log แบบ offline ไม่ใช่ decision authority

\subsection*{ถาม: ทำไมไม่ใช้ BO เหมือนงานเดิม?}
BO เหมาะกับการหา parameter ของ objective ที่ประเมินแพงแบบ offline ส่วนงานนี้มี discrete action set เพียง
16 mode vectors และต้อง enforce hard constraints ณ event จึงแจกแจงครบและ audit ได้ตรงกว่า

\subsection*{ถาม: วิธีนี้รับประกันลด transient หรือไม่?}
ยังไม่รับประกัน เป็น hypothesis ที่ต้องเทียบ peak frequency/RoCoF, voltage violation, current stress,
settling time และ handback jump กับ AGSI-only baseline บน input contract เดียวกัน

\subsection*{ถาม: IBR1 กลายเป็น slack bus หรือไม่?}
ไม่ IBR1 อาจเป็น reference owner ใน island แต่ยังเป็น dynamic current-injection device ใน full KCL
PF slack เป็นคนละความหมาย

\subsection*{ถาม: ถ้าไม่มี candidate ผ่านจะทำอย่างไร?}
ใช้ declared fail-closed fallback: hold current safe configuration หรือ verified emergency transaction
ถ้าผ่าน right-limit gates มิฉะนั้นคืน infeasible status ไม่ปลอม trajectory

\subsection*{ถาม: พร้อม real-time หรือยัง?}
ยัง ต้องมี runtime profiling และ deadline gate เพราะ short-horizon full-KCL prediction มีต้นทุนสูง

\section{ข้อจำกัดและ non-goals}

\begin{enumerate}
  \item ET-FCSPS ยังไม่ถูก implement และไม่มี closed-loop performance result
  \item ไม่เพิ่ม AVR; $E_{fd}$ และ voltage control contract ของ SG คงตาม model ปัจจุบัน
  \item ไม่แก้สมการ SG/IBR, AGSI++, threshold, case data หรือ event chronology ในขั้นรายงาน
  \item reserve/energy availability ต้องมี source/case definition ก่อนใช้เป็น production gate
  \item short horizon ไม่รับประกัน infinite-horizon stability โดยตัวมันเอง; ต้องมี falsification และ stability analysis เพิ่ม
  \item model mismatch, communication delay, protection hardware และ cyber-security อยู่นอกผลจำลองปัจจุบัน
  \item วิธีนี้ไม่ถูกกล่าวว่าดีกว่า BO ทุกกรณี และไม่ใช้ BO/LLM เป็น online switching authority
\end{enumerate}

\section{สรุป}

ET-FCSPS แก้ช่องว่างของ AGSI-only switching โดยแยก ``การตรวจพบความเครียดเฉพาะจุด'' ออกจาก
``การอนุมัติคำสั่งระดับระบบ'' AGSI++ ยังคงตรวจ local condition และ dwell ขณะที่ supervisor แจกแจง
mode/owner set, ตัด candidate ที่ผิด hard constraints, ทำนายผลช่วงสั้น แล้วเลือกคำสั่งที่ทำซ้ำและอธิบายได้
สำหรับสี่ IBR ขนาด candidate สูงสุด 16 mode vectors หรือ 32 mode-owner pairs ทำให้ exhaustive evaluation
เป็นทางเลือกที่ตรวจสอบได้โดยไม่ต้องใช้ black-box online optimiser

ประโยชน์ที่คาดหวังคือใช้ GFM เท่าที่จำเป็น ลด switching/transient/current stress ทำ handback นุ่มขึ้น
รักษา reference continuity และสร้าง evidence table ของทุกการตัดสินใจ แต่ทั้งหมดต้องพิสูจน์ด้วย implementation
และ paired comparison บน numerical contract เดียวกันก่อนจึงจะกล่าวเป็นผลของงานได้

\begin{thebibliography}{9}
\bibitem{rodriguez2013}
J. Rodriguez, M. P. Kazmierkowski, J. R. Espinoza, P. Zanchetta, H. Abu-Rub,
H. A. Young, and C. A. Rojas,
``State of the Art of Finite Control Set Model Predictive Control in Power Electronics,''
\emph{IEEE Transactions on Industrial Informatics}, vol. 9, no. 2, pp. 1003--1016, 2013.
DOI: 10.1109/TII.2012.2221469.

\bibitem{mayne2000}
D. Q. Mayne, J. B. Rawlings, C. V. Rao, and P. O. M. Scokaert,
``Constrained Model Predictive Control: Stability and Optimality,''
\emph{Automatica}, vol. 36, no. 6, pp. 789--814, 2000.
DOI: 10.1016/S0005-1098(99)00214-9.

\bibitem{frazier2018}
P. I. Frazier, ``A Tutorial on Bayesian Optimization,'' arXiv:1807.02811, 2018.

\bibitem{snoek2012}
J. Snoek, H. Larochelle, and R. P. Adams,
``Practical Bayesian Optimization of Machine Learning Algorithms,''
\emph{Advances in Neural Information Processing Systems 25}, 2012.

\bibitem{projectcode}
Power-flow repository: \path{+ibr/SwitchableIbr6.m},
\path{+stability/ibr_selector_table.m}, and the audited full-KCL hybrid runtime.
Accessed 4 August 2026.
\end{thebibliography}

\end{document}
